\section{问题重述}

\subsection{问题背景}

微网是为解决光伏发电等分布式能源的本地消纳、减少并网问题而发展起来的小型发用电系统。出于经济方面的考虑，非孤岛式微网可在外部电网（以下简称外网）电价较低或分布式能源有剩余电量的时段为储能设备充电，并在外网电价较高的时段释放电能；其调控目标是在利用小区负载、光伏发电量与储能当前储电量等数据的前提下，尽可能精准地给出微网的购电量与储能设备的充放电量，使微网提供的电力既满足小区负载、又尽可能节省购电费用。

题面附录给出储能设备参数：最大容量 12000 kWh，最大充放电功率 5000 kW，2025 年 1 月 1 日 0:00 的电量为 6000 kWh，运行过程中电量须保持在 $1200$--$10800$ kWh 之间，充放电效率为 90\%。

\subsection{问题提出}

\textbf{问题一}要求建立每天 0:00 制定当天计划购电策略的数学模型：电价与小区负载逐日相同、光伏发电功率有一天预测数据；微网提供的电能不得低于小区负载；储能设备在 0:00 与 24:00 的储电量相同；目标为全天购电费用最小。需基于附件 1 给出计划购电策略与充放电安排，保存为 \texttt{result1.xlsx}。

\textbf{问题二}把决策期扩展到全年：电价仍逐日相同，而负载与光伏逐日变化；若微网提供的电能低于负载，需向外网紧急购电，电价为交易时刻电价的 5 倍；除紧急购电外，其余购电费用按计划购电量计算。需基于附件 1 电价与附件 2 实际负载、实际光伏，在每天 0:00 制定当天计划，给出四个指定日期（2025 年 3 月 20 日、6 月 21 日、9 月 23 日、12 月 21 日）的结果与紧急购电结果，保存为 \texttt{result2.xlsx}。

\textbf{问题三}引入光伏功率的滚动预报：每天 0:00、6:00、12:00、18:00 可获得未来 24 小时整点预报。可依据 0:00 预报制定当天计划，也可依据其他时刻预报调整购电策略；计划购电量高于调整购电量时，违约电价为交易时刻电价的 50\%；调整购电量高于计划购电量时，超出部分电价为交易时刻电价的 1.5 倍；总费用包括计划购电费用、紧急购电费用与调整购电量的相关费用。需基于附件 1、2、3 建立模型、给出四个指定日期的结果并保存为 \texttt{result3.xlsx}，并分析是否需要引入其他时刻的预报来制定调整策略。

\textbf{问题四}指出外网电价本身也是实时波动的，要求在附件 4 的实时波动电价情景下，利用附件 2 的负载与光伏、附件 3 的光伏预报，\emph{重新计算问题二与问题三}，并分别保存为 \texttt{result4-2.xlsx}（对应问题二）与 \texttt{result4-3.xlsx}（对应问题三）。

由题面可见，问题四在数学上等价于\textbf{两道题}。二者共用附件 4 电价，但决策结构与结算方式不同：重算问题二只有每天一次的日计划与 5 倍价紧急购电，重算问题三则另有三个调整时刻与 $0.5\times/1.5\times$ 分段结算。因此本文自始至终按两条独立的结果链组织计算与验收，分别记为链 \emph{4-2} 与链 \emph{4-3}；两链共用同一价格预测器与同一结算方式（结算一律用实际电价）。
